\section{Introduction}
\label{sec:intro}
Evaluating a language model on a benchmark means running it on every item, and the cost multiplies when many checkpoints, benchmarks, prompts or inference settings are compared.
A natural shortcut is to run the model on a small subset and predict its full-benchmark score. Historical results of other (\emph{source}) models on every item are usually available and could guide which items to keep and how to extrapolate from them.
Several methods build on this idea: clustering items into anchors~\citep{vivek2024anchor}, IRT-based prediction~\citep{polo2024tiny,kipnis2025metabench}, learned predictors and selection policies~\citep{pacchiardi2024100,li2025active}. They were however evaluated on different model pools, splits, metrics and numbers of repetitions, often with few repetitions.

We ask how much the source models' results actually help, and when. Using the 0/1 item-level results of 13 benchmarks (Open LLM Leaderboard v2 and HELM), we run one protocol for all methods: the same budgets, the same source/target splits, selection from source models only, exact duplicate models removed, and uncertainty from model-level bootstrap with paired comparisons. Besides random splits of models (interpolation) we hold out the \emph{strongest} 20\% of models as targets (extrapolation) and, for the Open LLM benchmarks, whole organisations.
We re-implement the methods ourselves, and add a model-assisted sampling estimator, \mass{}, that adapts the difference estimator of survey sampling~\citep{sarndal1992model} to the benchmark setting: a stratified, probability-proportional-to-size item design built from source models, a latent-trait predictor, and a design-weighted correction on the sampled items.

\paragraph{Contributions and what we find.}
(1) A common, stricter protocol and re-implementations of the published methods we could reproduce and of simple baselines (Appendix~\ref{app:methods}); the data-repository conclusions of \citet{zhang2025benchmark} are tested as hypotheses and are confirmed or qualified in Section~\ref{sec:results}.
(2) \mass{}-Ens, whose design is a straightforward combination of known survey-sampling and prediction-powered ideas~\citep{sarndal1992model,angelopoulos2023ppi}; its value is empirical. With 50--100 target items it lowers the RMSE of the full-score prediction by about 21--27\% relative to the random-subset mean under extrapolation and interpolation, on all 13 benchmarks, an effective-sample-size gain of about 1.6--1.8. 
(3) Limits of that result: it is not distinguishable from the best baseline chosen by cross-validation on source models, nor from ridge regression on a random subset under interpolation, where ridge has the better point estimate; the lower error does not improve the ranking of target models; under extrapolation it is a small effect compared with the spread of the strongest models; and its reported intervals under-cover at small budgets.
(4) A decomposition showing that under extrapolation the difference correction on the sampled items is what prevents the degradation of predictors fitted on weaker models, with design and predictor each adding a little.
We conclude that subset-based prediction is effective under conditions: useful for estimating a score to a few accuracy points when the new model resembles the source models or a correction term is used, and not a reliable tool for ranking the strongest models.
